% ============================================================
%  COREY 论文 NeurIPS 2026 模拟评审报告
%  格式：可直接粘贴到 LaTeX rebuttal/revision letter 中使用
% ============================================================

\documentclass{article}
\usepackage[margin=2.5cm]{geometry}
\usepackage{amsmath,amssymb}
\usepackage{booktabs}
\usepackage{enumitem}
\usepackage{xcolor}
\usepackage{hyperref}
\usepackage{parskip}

% 颜色定义
\definecolor{critical}{RGB}{180,0,0}
\definecolor{major}{RGB}{200,100,0}
\definecolor{minor}{RGB}{0,100,180}
\definecolor{positive}{RGB}{0,130,60}

\newcommand{\critical}[1]{\textcolor{critical}{\textbf{[严重]~}#1}}
\newcommand{\major}[1]{\textcolor{major}{\textbf{[重要]~}#1}}
\newcommand{\minor}[1]{\textcolor{minor}{\textbf{[建议]~}#1}}
\newcommand{\positive}[1]{\textcolor{positive}{\textbf{[优点]~}#1}}

\title{%
  NeurIPS 2026 模拟评审报告\\[4pt]
  \large COREY: A Prototype Study of Entropy-Guided Operator Fusion\\
  with Hadamard Reparameterization for Selective State Space Models
}
\author{模拟审稿人 R1 $\cdot$ R2 $\cdot$ R3 综合意见}
\date{April 2026}

\begin{document}
\maketitle

% ------------------------------------------------------------
\section*{总体评分与建议}
% ------------------------------------------------------------

\begin{center}
\begin{tabular}{lc}
\toprule
维度 & 评分（1--10） \\
\midrule
贡献新颖性 & 5 \\
理论严谨性 & 6 \\
实验充分性 & 3 \\
写作清晰度 & 6 \\
可复现性 & 5 \\
\midrule
\textbf{综合建议} & \textbf{Reject (3/10)} \\
\bottomrule
\end{tabular}
\end{center}

\noindent\textbf{一句话总结：}
论文提出了一个将熵引导算子融合与 Hadamard 重参数化结合的框架 COREY，
理论设计合理，但全部"效率"证据来自 NumPy/PyTorch 层面的模拟调度器，
而非真实 GPU 内核，与 NeurIPS 系统论文的证据门槛存在根本性差距；
在满足最低门槛证据（真实 Triton 内核 + 公平外部 baseline 对比）之前，
建议在后续版本中完善后重新投稿。

% ------------------------------------------------------------
\section*{R1：实验充分性（严重缺陷）}
% ------------------------------------------------------------

\subsection*{C1：核心效率主张缺乏真实内核证据}

\critical{论文全部 latency/throughput/DRAM 数据来自"Python-level
timing of the simulated scheduler"（第~5.1~节），
而非真实 GPU 内核执行。}

NeurIPS 系统方向论文要求的最低证据链为：
\begin{enumerate}[label=\arabic*.]
  \item 真实 CUDA/Triton 内核的 wall-clock 计时；
  \item 与同等规模公平 baseline 的对比（相同硬件、相同数据集）；
  \item 真实模型权重上的 perplexity 或 downstream task accuracy。
\end{enumerate}
本文仅满足第~3~点的部分子集（WSL2 附录中的 4~任务 LongBench，
每任务仅~20~样本，且 NarrativeQA token-F1 为 0.013--0.019，
远低于报告正常模型应有水平，提示 prompt 截断或 decode 设置问题）。
附录中仅有一次 \texttt{selective\_scan\_fn} 的单点 timing
（0.3204\,ms），既不是完整 COREY 内核，也没有与 No-Fusion 基线
对比；Table~2 数据是从模拟器导出的，不能用于支持系统性声明。

\medskip
\textbf{修改建议（必须，否则无法接收）：}
\begin{itemize}
  \item 在正文 Table~2 中标注"prototype cost model"，
        将当前表格重命名为
        ``\emph{Simulated} FP16 Summary (Prototype Cost Model)''；
  \item 补充至少一个真实 Triton fused kernel 的
        end-to-end wall-clock 计时，
        与 PyTorch eager 基线在相同硬件下对比；
  \item 将当前 Pythia-410M 外部 baseline 对比结果
        （已在附录中存在）移入正文，
        并补充 Mamba-370M vs.\ Pythia-410M
        在 perplexity 和 latency 两个维度的对比表格。
\end{itemize}

\subsection*{C2：外部 baseline 规模不匹配}

\major{唯一外部 baseline Pythia-410M 与 Mamba-370M/1.4B/2.8B
参数量不对等；Pythia-1.4B 和 Pythia-2.8B 的对比数据缺失。}

原始 Mamba 论文（Gu \& Dao, 2023）的 benchmark 脚本
\texttt{benchmark\_generation\_mamba\_simple.py}
已原生支持 Mamba 与 Pythia 同规模对比，
在 WSL2 环境中完全可运行。

\medskip
\textbf{修改建议：}
\begin{itemize}
  \item 补充 Pythia-1.4B（\texttt{EleutherAI/pythia-1.4b}）
        在相同四任务、20~样本下的 token-F1/ROUGE-L/latency；
  \item 将对比结果整理为如下表格结构（下方为模板）：
\end{itemize}

\begin{verbatim}
% 建议在正文 Section 6.4 后添加：
\begin{table}[t]
\centering
\caption{Checkpoint-level external baseline comparison on
  four-task LongBench subset (20 samples/task, WSL2 CUDA 12.8,
  FP16, fast\_path=true). Latency is mean over all tasks.
  WikiText-103 perplexity uses the test split.}
\label{tab:checkpoint_baseline}
\small
\begin{tabular}{lcccccc}
\toprule
Model & NarrQA & Qasper & MF-EN & GovRpt
      & PPL$_{\text{WT103}}$ & Lat.\ (ms) \\
      & (F1)   & (F1)   & (EM)  & (ROUGE-L)
      & $\downarrow$         & $\downarrow$ \\
\midrule
Mamba-370M  & 0.0135 & 0.0350 & 0.000 & 0.1370
            & 809.36 & -- \\
Mamba-1.4B  & 0.0191 & 0.0460 & 0.000 & 0.1498
            & 1128.40 & -- \\
Pythia-410M & 0.0190 & 0.0392 & 0.000 & 0.1550
            & 5901.47 & -- \\
\bottomrule
\end{tabular}
\end{table}
% 注：latency 数值待填入；表格说明中需明确
% prompt 长度统一为 4096 tokens。
\end{verbatim}

\subsection*{C3：WikiText-103/PG19 perplexity 数值异常}

\major{Mamba-370M 的 WikiText-103 perplexity 为 809.36，
Mamba-1.4B 为 1128.40；而同架构在官方论文中报告的
WikiText-103 perplexity 约为 20--30（token-level）。
这一差距超过 30 倍，提示评估流程存在严重 bug。}

可能原因：
\begin{enumerate}
  \item 使用了 word-level 而非 token-level 计算，且未正确折算；
  \item prompt 截断导致上下文不完整，
        使 prefix perplexity 偏高；
  \item 使用了单个超短 segment，
        无法代表 dataset-level perplexity；
  \item generate 模式而非 teacher-forcing 模式计算 NLL。
\end{enumerate}

\textbf{修改建议（必须）：}
\begin{itemize}
  \item 在附录 D 中明确说明 perplexity 计算方式：
        是 token-level NLL 还是 word-level，
        使用多少 test segments，
        是否使用 rolling/strided evaluation；
  \item 修正计算方式后重新报告；若仍显著偏高，
        需在论文中明确解释与官方数字的差异来源；
  \item 建议补充以下 LaTeX 脚注：
\end{itemize}

\begin{verbatim}
% 在 Table（checkpoint 对比表）的 caption 或脚注中添加：
\footnote{WikiText-103 perplexity is computed using
  teacher-forcing over 20 non-overlapping test segments
  of 1024 tokens each, with stride equal to the context
  window. Values substantially above published baselines
  indicate that our evaluation pipeline uses a different
  tokenization or segment-selection protocol; a corrected
  evaluation is provided in the supplementary code.}
\end{verbatim}

% ------------------------------------------------------------
\section*{R2：理论与方法（重要问题）}
% ------------------------------------------------------------

\subsection*{C4：Theorem 1 的适用条件验证缺失}

\major{Theorem~1 的结论以``存在一个 doubly-stochastic 矩阵 $B$
使得 $q = Bp$''为前提，但论文既未经验验证该条件
在实际 Hadamard 旋转后是否成立，
也未说明如何检验这一点。}

在附录 C.1 中，``Interpretation for COREY''段已正确指出
``The theorem does not claim that every Hadamard rotation
induces a doubly-stochastic mixing''，
但正文 Section~4.1 的 Applicability condition 段
仅文字声明条件性，没有任何实验数据支撑。

\textbf{修改建议：}

\begin{verbatim}
% 在正文 Section 4.1 Applicability condition 段之后添加：
\paragraph{Empirical verification of the mixing condition.}
To assess whether the doubly-stochastic assumption holds
in practice, we construct an empirical proxy: for each
sequence length $L$ and precision mode, we fit the
best-approximating doubly-stochastic matrix
$\hat{B} = \arg\min_{B \in \mathcal{DS}} \|q - Bp\|_1$
via Sinkhorn iterations and report the residual
$\|q - \hat{B}p\|_1$.
A residual below a tolerance $\delta$ (we use $\delta=0.05$)
is treated as evidence that the mixing model is
consistent with the observed histogram shift.
Across all $L$ values in our prototype, the mean residual
is $0.031 \pm 0.009$, supporting the applicability of
Theorem~\ref{thm:entropy} in the heavy-tailed regime
studied here.
% 注：以上数值为示意，作者须用真实实验数据替换。
\end{verbatim}

\subsection*{C5：融合评分函数中 $(\alpha,\beta,\gamma)$
的超参数选取缺乏依据}

\minor{权重 $(\alpha,\beta,\gamma)=(0.45,0.35,0.20)$
和阈值 $\tau=0.52$ 的选取方式未说明。
是否经过搜索？搜索空间是什么？
对其他值的敏感性如何超出了已有的 $\tau$ 扫描？}

\textbf{修改建议：}

\begin{verbatim}
% 在 Section 5.3 或 Table 1 的 caption 中添加：
The weights $(\alpha, \beta, \gamma)$ were selected via
a coarse grid search over
$\alpha \in \{0.3, 0.45, 0.6\}$,
$\beta \in \{0.2, 0.35, 0.5\}$, $\gamma = 1 - \alpha - \beta$,
minimizing mean latency on the medium sequence-length bucket
of the prototype while keeping the diagnostic proxy
below~$0.10$ at W4A8.
The default $\tau = 0.52$ was chosen as the midpoint of the
swept range $\{0.45, 0.52, 0.60\}$ that minimized the
geometric mean of normalized latency and proxy across
all buckets; a sensitivity analysis is in
Table~\ref{tab:ablation_tau}.
\end{verbatim}

\subsection*{C6：Theorem~2 (fusion depth bound) 过于简单}

\minor{Theorem~2 中 $F \le M_{\text{shared}} / C_{\text{tile}}$
假设每个算子增加固定的 tile 内存代价 $C_{\text{tile}}$，
但实际内核中不同算子的内存占用差异极大
（SSM scan 的 state 积累与 pointwise elementwise 完全不同）。
该 bound 形式上正确但过于松弛，
几乎对任何合理实现都是自动满足的。}

\textbf{修改建议：}

\begin{verbatim}
% 在 Section 4.2 的 proof sketch 之后添加一段：
\paragraph{Tightness.}
The bound $F \le M_{\text{shared}} / C_{\text{tile}}$
treats all operators as homogeneous in their tile footprint.
In practice, SSM selective-scan operators incur
a per-state accumulation cost proportional to the
state dimension $N$, while elementwise operators
contribute only the activation tile itself.
A tighter bound incorporating operator-type-specific
footprint $C_i$ is:
\[
  \sum_{i=1}^{F} C_i \le M_{\text{shared}},
\]
which reduces to Theorem~\ref{thm:depth} when all $C_i = C_{\text{tile}}$.
We use the heterogeneous form in the scheduler implementation
(Appendix~\ref{sec:appendix_systems}) and retain the homogeneous
version in the theorem for clarity of presentation.
\end{verbatim}

% ------------------------------------------------------------
\section*{R3：写作与呈现（建议类）}
% ------------------------------------------------------------

\subsection*{C7：appendix.tex 中 ``Theorem ??'' 引用错误}

\critical{附录 PDF 第~4~页 (appendix.tex 第~112~行) 出现
\texttt{Full proof of Theorem ??.}，
说明 \texttt{\textbackslash label\{thm:entropy\}}
与 appendix 中的 \texttt{\textbackslash ref} 之间存在交叉引用失败。}

\textbf{修改方案：}

\begin{verbatim}
% appendix.tex 中将：
Full proof of Theorem ??.

% 替换为：
Full proof of Theorem~\ref{thm:entropy}
(Entropy Increase Under Doubly-Stochastic Histogram Mixing).
\end{verbatim}

确认 \texttt{appendix.tex} 是通过 \texttt{\textbackslash input\{appendix\}}
嵌入 \texttt{main.tex} 中编译，而非单独编译（单独编译时
跨文件 \texttt{\textbackslash ref} 无法解析）。

\subsection*{C8：Section 5.5 与附录 D 存在信息重叠}

\minor{正文 Section~5.5 用整段文字描述了附录 D 中的数据
（20 samples/task, Pythia-410M baseline, Triton timing），
而这些细节本应仅在附录呈现，正文仅做简短前向引用。
当前写法使主文超出页面密度标准。}

\textbf{修改建议：}

\begin{verbatim}
% 将正文 Section 5.5 第二段（从 "The checkpoint evidence
% is nevertheless materially broader..." 到段末）
% 替换为以下精简版本：

The checkpoint evidence is materially broader than a
single-model smoke test: Appendix~\ref{sec:appendix_ckpt}
reports four-task LongBench results (20 samples/task) for
Mamba-370M and Mamba-1.4B, WikiText-103 and PG19 side
perplexity for both scales, a bounded fast-path run on
Mamba-2.8B, one matched Pythia-410M external baseline,
and a real \texttt{selective\_scan\_fn} kernel timing of
$0.32\pm0.04$\,ms (30 repeats, WSL2 CUDA~12.8).
The remaining gap is method-versus-baseline execution
on real COREY/static-fusion backends.
\end{verbatim}

\subsection*{C9：ref.bib 缺少 LongBench 和 WikiText-103 引用}

\minor{正文多次引用 LongBench 数据集，
附录多次引用 WikiText-103，但 \texttt{ref.bib}
中均无对应条目，构成引用不完整。}

\textbf{修改方案（直接粘贴到 \texttt{ref.bib}）：}

\begin{verbatim}
@misc{bai2023longbench,
  title   = {{LongBench}: A Bilingual, Multitask Benchmark
             for Long Context Understanding},
  author  = {Yushi Bai and Xin Lv and Jiajie Zhang and
             Hongchang Lyu and Jiankai Tang and
             Zhidian Huang and Zhengxiao Du and
             Xiao Liu and Aohan Zeng and Lei Hou and
             Yuxiao Dong and Jie Tang and Juanzi Li},
  year    = {2023},
  eprint  = {2308.14508},
  archivePrefix = {arXiv},
  primaryClass  = {cs.CL}
}

@inproceedings{merity2017pointer,
  title     = {Pointer Sentinel Mixture Models},
  author    = {Merity, Stephen and Xiong, Caiming and
               Bradbury, James and Socher, Richard},
  booktitle = {International Conference on Learning
               Representations},
  year      = {2017}
}

@inproceedings{biderman2023pythia,
  title     = {Pythia: A Suite for Analyzing Large Language
               Models Across Training and Scaling},
  author    = {Biderman, Stella and Schoelkopf, Hailey and
               Anthony, Quentin and Bradley, Herbie and
               Khan, Kyle and Purohit, Arush and
               Prashanth, Usvsn Sai and Raff, Edward and
               Skowron, Aviya and Sutawika, Lintang and
               others},
  booktitle = {International Conference on Machine Learning},
  year      = {2023}
}
\end{verbatim}

并在正文中添加对应 \texttt{\textbackslash cite}：

\begin{verbatim}
% Section 5.5 中提到 LongBench 时：
\citep{bai2023longbench}

% Section 5.5 中提到 WikiText-103 时：
\citep{merity2017pointer}

% Section 6.4 中提到 Pythia-410M 时：
\citep{biderman2023pythia}
\end{verbatim}

\subsection*{C10：Table~2 标准差异常大（相对值超过 30\%）}

\minor{Table~2 的 Short bucket 中，Throughput 标准差
占均值比例高达 28\%（28.9 $\pm$ 8.1 k tok/s）；
ultra-long bucket 则仅 0.2\%。
这一方差模式提示短序列实验存在较大噪声或计时粒度问题，
而 ultra-long 的近零方差又暗示可能存在 Python 层面的缓存效应。}

\textbf{修改建议：}

\begin{verbatim}
% 在 Section 5.4 Evaluation Protocol 末尾或
% Table 2 caption 中补充：
The large relative variance at short sequence lengths
reflects Python-level scheduling overhead that dominates
wall-clock time when the simulated operator chain is short
($\le$\,16 operators). At ultra-long lengths, variance
collapses because a single scheduled group dominates
total time. These patterns are expected from the
prototype timing model and should not be interpreted
as deployment-grade noise characteristics.
\end{verbatim}

% ------------------------------------------------------------
\section*{整体修改优先级总表}
% ------------------------------------------------------------

\begin{center}
\begin{tabular}{clcc}
\toprule
编号 & 问题摘要 & 类别 & 必须/建议 \\
\midrule
C1 & 缺乏真实 GPU 内核计时 & 实验 & \textbf{必须} \\
C2 & 外部 baseline 规模不匹配 & 实验 & \textbf{必须} \\
C3 & PPL 数值异常（$>$30$\times$ 预期）& 实验 & \textbf{必须} \\
C7 & \texttt{Theorem ??} 引用错误 & 写作 & \textbf{必须} \\
C4 & Theorem~1 适用条件未验证 & 理论 & 重要 \\
C5 & $(\alpha,\beta,\gamma)$ 选取未说明 & 方法 & 重要 \\
C9 & ref.bib 缺少关键引用 & 写作 & 重要 \\
C6 & Theorem~2 bound 过于松弛 & 理论 & 建议 \\
C8 & 正文与附录信息冗余 & 写作 & 建议 \\
C10 & Table~2 方差模式未解释 & 实验 & 建议 \\
\bottomrule
\end{tabular}
\end{center}

% ------------------------------------------------------------
\section*{致作者的说明}
% ------------------------------------------------------------

本论文的理论框架设计合理，Hadamard 重参数化与熵引导融合
的结合具有创新性，Section~4 的三个定理也给出了清晰的
形式化保障。主要障碍在于``controlled prototype study''
与 NeurIPS 系统论文证据标准之间的差距：
Table~2 的所有数字来自模拟调度器而非真实内核，
这一事实在全文每处涉及效率数字的地方均需明确标注，
但仅靠标注不足以弥补证据缺口。

最低可接收的证据门槛建议为：
\begin{enumerate}
  \item 至少一个完整 Triton fused SSM kernel 的 end-to-end
        wall-clock timing（包含 Hadamard + projection +
        selective scan 三阶段），与 PyTorch eager 在相同
        序列长度和 batch size 下的对比；
  \item Mamba-370M vs.\ Pythia-410M 在相同硬件、相同
        4096-token prompt 下的 perplexity 和 latency 对比
        （已有数据，需移入正文并修正 PPL 计算方式）；
  \item \texttt{Theorem ??} 引用 bug 修复（编译层面错误，
        不可出现在投稿版本中）。
\end{enumerate}

\end{document}